\scanpage{60}
% @01–05, hypercohomology
Idem \deleted{pour} $A_*$ muni d'une structure suppl -- p ex. de groupe abélien, éventuellement diff et gradué. Dans ce dernier cas -- si les degrés sont bornés \unsure{infért} -- on peut donc prendre l'hypercohomologie.
% @06, definition
\[
 \mathbb{R}^*\unsure{\mathbb{\Gamma}|\Gamma}_{\unsure{?|,}}(\unsure{\mathscr{X}|\mathcal{X}|X},A_*^*)\overset{\mathrm{dfn}}{=}
 \mathbb{R}^*\Gamma_{\mathscr{X}}(\underline{A_{\mathscr{X}}^*})
\]
% @07–09, gloss below right term
\margin{complexe induit par $A_*^*$ sur $\mathscr{X}$ ou $A_*^*\times\mathscr{X}/\mathscr{X}$}
% @10–11, cochains
Soit $A$\inserted{$:(\deleted{\unsure{\Delta}}\mathbb{\Delta}/\mathscr{X})^{\circ}\longrightarrow(\mathrm{Ab})$} un groupe sur $\mathscr{X}$, considérons le complexe des cochaines sur $\mathscr{X}$ à coeff. ds $A$
% @12–13, product
\[
 C^n(\mathscr{X};A)\overset{\mathrm{dfn}}{=}
 \prod_{\substack{\sigma\in\operatorname{Ob}\,\Delta_n/\mathscr{X}\\=\mathscr{X}_n}}A(\sigma),
\]
% @14–15, final observation
Les $C^n$ forment un groupe co $1/2$ simplicial
% archive sheet number: 59

